\documentclass[11pt,a4paper]{article}
\usepackage[utf8]{inputenc}\usepackage[T1]{fontenc}\usepackage{textcomp}\usepackage[french]{babel}
\usepackage{lmodern}\usepackage[top=2cm,bottom=2cm,left=2.2cm,right=2.1cm]{geometry}
\usepackage{microtype}\usepackage{xcolor}\usepackage{enumitem}\usepackage{titlesec}
\usepackage[most]{tcolorbox}\usepackage{hyperref}
\definecolor{mainblue}{RGB}{28,74,102}\definecolor{darkgray}{RGB}{62,62,66}
\definecolor{ent}{RGB}{221,236,247}\definecolor{ext}{RGB}{250,233,207}\definecolor{mod}{RGB}{238,238,242}
\titleformat{\section}{\large\bfseries\color{mainblue}}{\thesection}{0.6em}{}
\setlist[itemize]{leftmargin=1.2em,itemsep=0.05em,topsep=0.1em}
\setlength{\parindent}{0pt}\setlength{\parskip}{0.3em}
\newtcolorbox{client}{colback=ext!60,colframe=ext!60!black,boxrule=0.4pt,arc=1.5mm,left=2mm,right=2mm,
  top=0.8mm,bottom=0.8mm,width=0.78\textwidth,enlarge left by=0.22\textwidth,
  fontupper=\small,title={Client},fonttitle=\scriptsize\bfseries,coltitle=darkgray,colbacktitle=ext!60}
\newtcolorbox{assistant}{colback=ent,colframe=mainblue!60,boxrule=0.4pt,arc=1.5mm,left=2mm,right=2mm,
  top=0.8mm,bottom=0.8mm,width=0.85\textwidth,fontupper=\small,
  title={Assistant},fonttitle=\scriptsize\bfseries,coltitle=white,colbacktitle=mainblue!80}
\newtcolorbox{pourquoi}{colback=mod,colframe=mod,boxrule=0pt,arc=1mm,left=2mm,right=2mm,top=0.5mm,
  bottom=0.5mm,fontupper=\scriptsize,title={Pourquoi},fonttitle=\scriptsize\bfseries,coltitle=darkgray,
  colbacktitle=mod}
\newcommand{\frugal}[1]{{\scriptsize\color{darkgray}\textbf{Frugalité} : #1}\par\medskip}
\begin{document}
{\LARGE\bfseries\color{mainblue} Démo de secours\par}
{\large Assistant bancaire frugal --- le script de la démonstration, rejoué par le vrai pipeline\par}
{\small\color{darkgray} Généré le 27/10/2026 à 10 h 16 sur yahya-bel-hajjam (x86\_64), mode : gabarits, sans LLM.\par}
\medskip
\section*{Bilan}
15 tours de conversation, dont \textbf{15 sans aucun appel au LLM} ; 38,0~ms de calcul au total, soit 2,53~ms par tour en moyenne ; 0 appel(s) au LLM ; 0 escalade. Mode : gabarits, sans LLM.
\section{Carte perdue : réagir vite et agir}
\textit{Inès, 34 ans, cadre (C01).} Compris par les mots-clés ; l'opposition n'est faite qu'après le « oui ».

\begin{client}J'ai perdu ma carte hier soir, je fais quoi ?\end{client}
\begin{assistant}Je comprends, on sécurise votre carte tout de suite. Je peux faire opposition sur votre carte maintenant : plus aucun paiement ne sera possible avec. Confirmez-vous ? (oui / non)\end{assistant}
\begin{pourquoi}
\begin{itemize}
\item Intention \texttt{lost\_or\_stolen\_card}, comprise par mots\_cles
\item Action \texttt{bloquer\_carte} : en attente du « oui »
\end{itemize}
\end{pourquoi}
\frugal{comprendre 0,3~ms, converser 26,4~ms, formuler 0,0~ms --- appels au LLM : 0}

\begin{client}oui\end{client}
\begin{assistant}C'est fait : votre carte est en opposition depuis 10 h 16 (référence OPP-C01-0001). Plus personne ne peut s'en servir, et les opérations frauduleuses faites après l'opposition ne sont pas à votre charge. Une nouvelle carte peut être commandée depuis votre espace client ou avec un conseiller. Rappel : la banque ne vous demandera jamais votre code confidentiel.\end{assistant}
\begin{pourquoi}
\begin{itemize}
\item Action \texttt{bloquer\_carte} : exécutée
\end{itemize}
\end{pourquoi}
\frugal{agir 0,1~ms, formuler 0,0~ms --- appels au LLM : 0}

\section{L'appartement : questionner, calculer, expliquer}
\textit{Léa, 31 ans, cadre (C07).} Les revenus ne sont pas redemandés ; refus chiffré, puis correction recalculée.

\begin{client}Je veux acheter un appartement\end{client}
\begin{assistant}C'est un beau projet. J'ai noté vos revenus de 2\,900~\texteuro{} nets par mois. Quel est le prix du bien que vous visez, et quelle part de vos 25\,000~\texteuro{} d'épargne souhaitez-vous y consacrer ?\end{assistant}
\begin{pourquoi}
\begin{itemize}
\item Intention \texttt{credit\_immobilier}, comprise par mots\_cles
\item Information demandée : \texttt{prix\_bien}
\end{itemize}
\end{pourquoi}
\frugal{comprendre 0,2~ms, converser 0,1~ms, formuler 0,0~ms --- appels au LLM : 0}

\begin{client}Autour de 250 000, et je peux mettre les 25 000\end{client}
\begin{assistant}Avec un prix de 250\,000~\texteuro{} et un apport de 25\,000~\texteuro{}, il faudrait emprunter environ 243\,750~\texteuro{}, frais de notaire estimés à 7,5~\% compris. Sur 25 ans, à un taux indicatif de 3,5~\%, la mensualité serait d'environ 1\,220~\texteuro{}, soit 42,1~\% de vos revenus : c'est au-delà des 35~\% retenus pour un crédit immobilier. Avec votre apport, un bien autour de 212\,000~\texteuro{} resterait dans cette limite. Un conseiller peut étudier les leviers : un apport plus élevé, un co-emprunteur, des aides comme le PTZ selon conditions. Voulez-vous qu'un conseiller vous rappelle ? (oui / non) Un crédit vous engage et doit être remboursé. Vérifiez vos capacités de remboursement avant de vous engager.\end{assistant}
\begin{pourquoi}
\begin{itemize}
\item Intention \texttt{credit\_immobilier}, comprise par reponse\_attendue
\item Proposé : Rendez-vous avec un conseiller --- Avec votre apport, un bien autour de 212\,000~\texteuro{} resterait sous les 35~\%. Un conseiller étudie les leviers (apport, co-emprunteur, aides comme le PTZ selon conditions) et décide.
\item Écarté : Crédit immobilier --- mensualité d'environ 1\,220~\texteuro{}, soit 42,1~\% de vos revenus (maximum 35~\%) (\texttt{HCSF-taux-effort})
\item Écarté : Un prêt immobilier sur 30 ans --- au-delà de la durée maximale de 25 ans (\texttt{HCSF-duree})
\item Alerte : Taux d'effort de 42,1~\% : au-delà des 35~\% retenus pour un crédit immobilier.
\end{itemize}
\end{pourquoi}
\frugal{comprendre 0,8~ms, converser 0,1~ms, conseiller 0,4~ms, formuler 0,1~ms --- appels au LLM : 0}

\begin{client}Et avec un bien à 200 000 ?\end{client}
\begin{assistant}Avec un prix de 200\,000~\texteuro{} et un apport de 25\,000~\texteuro{}, il faudrait emprunter environ 190\,000~\texteuro{}, frais de notaire estimés à 7,5~\% compris. Sur 25 ans, à un taux indicatif de 3,5~\%, la mensualité serait d'environ 951~\texteuro{}, soit 32,8~\% de vos revenus : c'est sous les 35~\%, le projet paraît finançable. Voulez-vous qu'un conseiller vous rappelle ? (oui / non) Un crédit vous engage et doit être remboursé. Vérifiez vos capacités de remboursement avant de vous engager.\end{assistant}
\begin{pourquoi}
\begin{itemize}
\item Intention \texttt{credit\_immobilier}, comprise par reponse\_attendue
\item Proposé : Crédit immobilier --- Mensualité estimée à 951~\texteuro{} sur 25 ans, soit 32,8~\% de vos revenus : sous les 35~\%.
\item Proposé : Rendez-vous avec un conseiller --- Avec votre apport, un bien autour de 212\,000~\texteuro{} resterait sous les 35~\%. Un conseiller étudie les leviers (apport, co-emprunteur, aides comme le PTZ selon conditions) et décide.
\item Écarté : Un prêt immobilier sur 30 ans --- au-delà de la durée maximale de 25 ans (\texttt{HCSF-duree})
\end{itemize}
\end{pourquoi}
\frugal{comprendre 0,2~ms, converser 0,1~ms, conseiller 0,1~ms, formuler 0,0~ms --- appels au LLM : 0}

\section{Même phrase, deux clients : Sophie}
\textit{Sophie, 41 ans, employée administrative (C03).} Pour ses enfants : Livret A, Livret Jeune pour l'aîné ; LDDS et LEP écartés.

\begin{client}Je voudrais mettre de côté pour mes enfants.\end{client}
\begin{assistant}Bonne idée de mettre de côté. Dans combien de temps pensez-vous avoir besoin de cet argent ?\end{assistant}
\begin{pourquoi}
\begin{itemize}
\item Intention \texttt{epargne\_projet}, comprise par mots\_cles
\item Information demandée : \texttt{horizon\_mois}
\end{itemize}
\end{pourquoi}
\frugal{comprendre 0,2~ms, converser 0,1~ms, formuler 0,0~ms --- appels au LLM : 0}

\begin{client}Dans 10 ans, 100 euros par mois\end{client}
\begin{assistant}Voici ce que je vous propose. Livret A : ouvrable au nom de chaque enfant mineur ; argent disponible ; intérêts sans impôt. Livret Jeune : votre enfant de 12 ans peut en ouvrir un (de 12 à 25 ans). Je ne vous propose pas : Livret de développement durable et solidaire (LDDS), réservé aux majeurs ; Livret d'épargne populaire (LEP), réservé aux personnes sous un plafond de revenu fiscal, pas ouvert à un enfant.\end{assistant}
\begin{pourquoi}
\begin{itemize}
\item Intention \texttt{epargne\_projet}, comprise par reponse\_attendue
\item Proposé : Livret A --- Ouvrable au nom de chaque enfant mineur ; argent disponible ; intérêts sans impôt.
\item Proposé : Livret Jeune --- Votre enfant de 12 ans peut en ouvrir un (de 12 à 25 ans).
\item Écarté : Livret de développement durable et solidaire (LDDS) --- réservé aux majeurs (\texttt{LDDS-age})
\item Écarté : Livret d'épargne populaire (LEP) --- réservé aux personnes sous un plafond de revenu fiscal, pas ouvert à un enfant (\texttt{LEP-revenu})
\end{itemize}
\end{pourquoi}
\frugal{comprendre 0,2~ms, converser 0,1~ms, conseiller 0,0~ms, formuler 0,1~ms --- appels au LLM : 0}

\section{Même phrase, deux clients : Lucas}
\textit{Lucas, 20 ans, étudiant (C08).} Pour lui, 20 ans : le Livret Jeune passe en premier.

\begin{client}Je voudrais mettre de côté\end{client}
\begin{assistant}Bonne idée de mettre de côté. Dans combien de temps pensez-vous avoir besoin de cet argent ?\end{assistant}
\begin{pourquoi}
\begin{itemize}
\item Intention \texttt{epargne\_projet}, comprise par mots\_cles
\item Information demandée : \texttt{horizon\_mois}
\end{itemize}
\end{pourquoi}
\frugal{comprendre 0,2~ms, converser 0,1~ms, formuler 0,0~ms --- appels au LLM : 0}

\begin{client}Dans 10 ans, 100 euros par mois\end{client}
\begin{assistant}Voici ce que je vous propose. Livret Jeune : vous avez 20 ans : le Livret Jeune est réservé aux 12-25 ans. Livret A : argent disponible à tout moment ; intérêts sans impôt. Livret de développement durable et solidaire (LDDS) : argent disponible à tout moment ; intérêts sans impôt. Je ne vous propose pas : Livret d'épargne populaire (LEP), éligibilité à vérifier : le plafond de revenu n'est pas encore renseigné ; Plan épargne logement (PEL), adapté à un projet immobilier ; l'argent n'est pas disponible avant la fin du plan.\end{assistant}
\begin{pourquoi}
\begin{itemize}
\item Intention \texttt{epargne\_projet}, comprise par reponse\_attendue
\item Proposé : Livret Jeune --- Vous avez 20 ans : le Livret Jeune est réservé aux 12-25 ans.
\item Proposé : Livret A --- Argent disponible à tout moment ; intérêts sans impôt.
\item Proposé : Livret de développement durable et solidaire (LDDS) --- Argent disponible à tout moment ; intérêts sans impôt.
\item Écarté : Livret d'épargne populaire (LEP) --- éligibilité à vérifier : le plafond de revenu n'est pas encore renseigné (\texttt{LEP-revenu})
\item Écarté : Plan épargne logement (PEL) --- adapté à un projet immobilier ; l'argent n'est pas disponible avant la fin du plan (\texttt{PEL-objectif})
\end{itemize}
\end{pourquoi}
\frugal{comprendre 0,2~ms, converser 0,1~ms, conseiller 0,0~ms, formuler 0,0~ms --- appels au LLM : 0}

\section{Les frais de Marc : expliquer, inclure}
\textit{Marc, 52 ans, ouvrier (C02).} Lu dans l'historique ; plafond légal ; offre clientèle fragile proposée.

\begin{client}Pourquoi j'ai payé 8 \texteuro{} de frais ce mois-ci ?\end{client}
\begin{assistant}Le 8 octobre, une commission d'intervention de 8~\texteuro{} : un paiement est passé alors que votre compte dépassait son découvert autorisé. La loi plafonne ces commissions à 8~\texteuro{} par opération et 80~\texteuro{} par mois. Vous avez eu 5 découverts en 12 mois : un conseiller peut vous présenter l'offre clientèle fragile, qui limite ces frais à 4~\texteuro{} par opération et 20~\texteuro{} par mois. Voulez-vous qu'il vous rappelle ? (oui / non)\end{assistant}
\begin{pourquoi}
\begin{itemize}
\item Intention \texttt{extra\_charge\_on\_statement}, comprise par mots\_cles
\item Action \texttt{expliquer\_frais} : exécutée
\end{itemize}
\end{pourquoi}
\frugal{comprendre 0,2~ms, converser 0,1~ms, agir 0,0~ms, formuler 0,0~ms --- appels au LLM : 0}

\section{Julie, mécontente : passer la main}
\textit{Julie, 37 ans, employée de services (C06).} Pas de justification : un conseiller reprend la main avec le résumé.

\begin{client}Votre appli ne marche jamais !\end{client}
\begin{assistant}Je suis désolé pour ce désagrément. C'est votre 3e message aujourd'hui : je préfère qu'un conseiller reprenne la main, avec le résumé de vos échanges. Voulez-vous qu'il vous rappelle ? (oui / non)\end{assistant}
\begin{pourquoi}
\begin{itemize}
\item Intention \texttt{insatisfaction}, comprise par mots\_cles
\item Action \texttt{transmettre\_conseiller} : en attente du « oui »
\end{itemize}
\end{pourquoi}
\frugal{comprendre 0,2~ms, converser 0,0~ms, formuler 0,0~ms --- appels au LLM : 0}

\begin{client}oui\end{client}
\begin{assistant}C'est noté : un conseiller va vous rappeler (demande TCK-C06-0002). Il aura le résumé de notre échange, vous n'aurez pas à tout répéter.\end{assistant}
\begin{pourquoi}
\begin{itemize}
\item Action \texttt{transmettre\_conseiller} : exécutée
\end{itemize}
\end{pourquoi}
\frugal{agir 0,0~ms, formuler 0,0~ms --- appels au LLM : 0}

\section{Paul, 68 ans : pas à pas, en mots simples}
\textit{Paul, 68 ans, retraité (C04).} Mode « mots simples » automatique ; rien ne part sans confirmation.

\begin{client}Je dois envoyer 500 \texteuro{} à mon fils.\end{client}
\begin{assistant}D'accord, je vous guide pas à pas.\par Je récapitule : un virement de 500~\texteuro{} vers votre fils, Thomas.\par Confirmez-vous ? (oui / non) Rappel : la banque ne vous demandera jamais votre code confidentiel.\end{assistant}
\begin{pourquoi}
\begin{itemize}
\item Intention \texttt{faire\_virement}, comprise par mots\_cles
\item Action \texttt{preparer\_virement} : en attente du « oui »
\end{itemize}
\end{pourquoi}
\frugal{comprendre 0,2~ms, converser 0,0~ms, formuler 1,1~ms --- appels au LLM : 0}

\begin{client}oui\end{client}
\begin{assistant}C'est fait : le virement de 500~\texteuro{} vers Thomas est enregistré, référence VIR-C04-0003.\end{assistant}
\begin{pourquoi}
\begin{itemize}
\item Action \texttt{preparer\_virement} : exécutée
\end{itemize}
\end{pourquoi}
\frugal{agir 0,0~ms, formuler 0,3~ms --- appels au LLM : 0}

\section{Le jargon, même mal écrit}
\textit{Lucas, 20 ans, étudiant (C08).} Fiche du glossaire retrouvée malgré la faute de frappe, sans LLM.

\begin{client}C'est quoi le TEAG ?\end{client}
\begin{assistant}TAEG : Le coût total d'un crédit (intérêts, frais, assurance obligatoire) exprimé en pourcentage par an. C'est le chiffre qui permet de comparer deux offres.\end{assistant}
\begin{pourquoi}
\begin{itemize}
\item Intention \texttt{question\_definition}, comprise par mots\_cles
\end{itemize}
\end{pourquoi}
\frugal{comprendre 0,2~ms, converser 0,0~ms, formuler 5,5~ms --- appels au LLM : 0}

\end{document}
